\chapter{Revisão da literatura}
\label{sec_Revisao_literatura}

Este capítulo apresenta a revisão da literatura relacionada aos principais temas envolvidos nesta dissertação, com foco nas operações de \textit{rendezvous} e \textit{berthing}, na dinâmica relativa entre espaçonaves, na modelagem e no controle de manipuladores robóticos espaciais, bem como nas estratégias de guiamento, planejamento e autonomia aplicadas a esse contexto.

Inicialmente, são apresentadas as obras de referência e os fundamentos conceituais que sustentam o desenvolvimento do tema. Em seguida, é feita uma revisão da literatura relacionada à evolução dos estudos sobre \textit{rendezvous} e aproximação relativa, bem como às técnicas de estimação de pose do alvo e percepção relativa em operações de proximidade. Posteriormente, são revisados trabalhos voltados à modelagem e ao controle de sistemas acoplados espaçonave--manipulador, incluindo contribuições relacionadas ao controle com reação nula em manipuladores \textit{free-floating}. Na sequência, são apresentadas referências mais recentes sobre planejamento, autonomia e técnicas aplicadas a operações de proximidade orbital. Ao final do capítulo, é apresentada uma síntese crítica da literatura revisada, com destaque para a lacuna identificada e para o posicionamento da presente dissertação.


\section{Fundamentos e obras de referência}
\label{sec_fundamentos_obras_referencia}

Nesta seção são apresentadas as principais obras de referência que fundamentam os conceitos, as formulações matemáticas e os procedimentos de modelagem utilizados ao longo desta dissertação.

No contexto das operações de proximidade orbital, a formulação clássica da dinâmica relativa é estabelecida por \cite{clohessy1960terminal}. 
Neste trabalho seminal, os autores desenvolveram as equações lineares do movimento relativo entre duas espaçonaves, conhecidas como equações de \gls{hcw}. 
A importância dessa formulação na literatura é amplamente reconhecida, pois ela fornece a base matemática essencial para análises preliminares, projetos de guiamento terminal e desenvolvimento de leis de controle em operações de aproximação.
Entretanto, a aplicação do modelo possui limitações estritas: a formulação assume que a órbita do alvo é perfeitamente circular, ou com excentricidade desprezível, e restringe-se a distâncias relativas curtas, nas quais a aproximação linear do campo gravitacional se mantém válida.
Mesmo com essas restrições, o modelo consolidou o tratamento matemático da dinâmica de \gls{rvdb} e serve de ponto de partida para todas as formulações com maior nível de detalhamento.

No que se refere à dinâmica e ao controle de atitude de espaçonaves, \cite{wie2008space}, \cite{hughes1986spacecraft} e \cite{wertz1978spacecraft} apresentam os fundamentos da dinâmica rotacional de corpos rígidos, as diferentes formas de representação da atitude e os principais conceitos empregados na modelagem e no controle de espaçonaves.
De forma particular, \cite{wie2008space} apresenta a formulação da dinâmica de atitude em ambiente orbital, incluindo a relação entre os referenciais inercial, orbital e fixo ao corpo. Além disso, essas referências também discutem os principais torques atuantes no sistema, a formulação das equações de Euler e os efeitos associados ao movimento orbital do referencial local.

As diferentes representações da atitude, tais como ângulos de Euler, cossenos diretores, quatérnios e parametrizações equivalentes, são discutidas por \cite{wertz1978spacecraft}, \cite{hughes1986spacecraft}, \cite{wie2008space} e \cite{junkins1986optimal}. Essas obras apresentam as vantagens e limitações de cada parametrização, bem como sua aplicação na formulação cinemática e dinâmica de problemas de atitude. Essas referências estabelecem o formalismo matemático empregado na descrição da orientação da espaçonave perseguidora ao longo deste trabalho.

No campo da robótica, \cite{craig2005} e \cite{sciavicco_siciliano2010} apresentam os fundamentos da modelagem cinemática e dinâmica de manipuladores robóticos, incluindo a descrição de elos e juntas, o uso de parâmetros geométricos, a obtenção de transformações homogêneas e a formulação das equações do movimento por diferentes abordagens.

As formulações voltadas especificamente à robótica espacial têm como marcos os trabalhos de \cite{umetani1989resolved}, \cite{papadopoulos1991nature} e \cite{dubowsky1993kinematics}. 
Diferentemente dos manipuladores fixos em solo, esses autores demonstram que os movimentos das juntas produzem reações sobre a base espacial, afetando a atitude e o movimento translacional. 
\cite{umetani1989resolved} introduziram a matriz Jacobiana generalizada para o controle cinemático em regime \textit{free-floating}, permitindo relacionar as velocidades articulares ao movimento do efetuador terminal mesmo quando a espaçonave reage dinamicamente. 
\cite{papadopoulos1991nature} aprofundaram a natureza dos algoritmos de controle para esses sistemas, mostrando que a ausência de base fixa introduz restrições singulares na formulação das leis de controle. 
Por fim, \cite{dubowsky1993kinematics} apresentaram a sistematização definitiva da modelagem cinemática e dinâmica, estabelecendo as diferenças conceituais entre sistemas \textit{free-flying} e \textit{free-floating}. 
Esses trabalhos consolidam a base conceitual necessária para a compreensão do acoplamento dinâmico entre a espaçonave e o manipulador robótico.

De modo geral, as obras de referência aqui apresentadas fornecem os fundamentos conceituais e matemáticos necessários para a análise da literatura especializada discutida nas seções seguintes. A partir dessas referências, torna-se possível examinar com maior clareza os trabalhos publicados sobre \gls{rvdb}, percepção relativa do alvo, modelagem acoplada de sistemas espaçonave--manipulador e estratégias de controle e autonomia aplicadas a operações orbitais de proximidade.
% livros-chave;
% bases conceituais;
% definições centrais;
% formulações clássicas que sustentam o restante.

\section{Estudos sobre \gls{rvdb} e aproximação relativa}
\label{sec_evolucao_rendezvous_aproximacao_relativa}

Historicamente, o estudo do movimento relativo entre espaçonaves e o planejamento de manobras de aproximação consolidaram-se a partir da formulação clássica de \cite{clohessy1960terminal}. 
No entanto, ao longo do desenvolvimento da área, observou-se que o problema de \gls{rvdb} não poderia permanecer restrito apenas à dinâmica relativa linearizada.
Em missões de serviço em órbita, inspeção, captura e acoplamento, passou a ser necessário considerar não apenas o movimento relativo geométrico entre perseguidor e alvo, mas também restrições operacionais, requisitos de segurança e diferentes níveis de cooperação do alvo durante a operação.

Nesse contexto, \cite{arantes2010far} apresentam um estudo sobre manobras \textit{far} e \textit{proximity} de uma constelação de satélites de serviço, tratando de forma integrada diferentes fases da missão. Os autores empregam a dinâmica relativa baseada nas equações de \gls{hcw} e discutem tanto as manobras distantes quanto as operações de proximidade. O trabalho incorpora ainda a estimação autônoma da pose do satélite cliente, mostrando que, em missões de serviço orbital, a aproximação não depende apenas da formulação da trajetória relativa, mas também da capacidade de acompanhar a configuração espacial do alvo. Os resultados apresentados pelos autores mostram a viabilidade da abordagem proposta para integrar análise orbital e operação de proximidade em um mesmo contexto de \gls{rvdb}.

Mais recentemente, \cite{fiori2024modeling} apresentam um arcabouço de modelagem, simulação e controle para \gls{rvdb} automatizado nas proximidades de uma estação espacial, considerando restrições posicionais e obstáculos físicos. Nesse trabalho, o movimento da espaçonave é formulado por meio de um formalismo geométrico baseado em grupos de Lie, enquanto o problema de aproximação é tratado com o uso de potenciais virtuais. Os autores realizam simulações numéricas com o objetivo de comparar estratégias de controle e avaliar seu desempenho no cenário estudado.

De modo geral, a literatura revisada mostra que os estudos sobre \gls{rvdb} e aproximação relativa evoluíram de formulações clássicas centradas na dinâmica relativa linearizada para abordagens mais amplas, nas quais passam a ser considerados aspectos operacionais, geométricos e autônomos da missão. Enquanto os trabalhos clássicos estabeleceram a base matemática para o tratamento do movimento relativo, estudos posteriores ampliaram o escopo do problema ao incorporar requisitos de segurança, interação com alvos cooperativos e não cooperativos, além de integração com tarefas de navegação e operação de proximidade. Esses aspectos são relevantes para a presente dissertação, pois mostram que a fase de aproximação curta deve ser analisada dentro de um contexto mais amplo, no qual dinâmica relativa, atitude, restrições operacionais e interação com o alvo tornam-se progressivamente acopladas.
% artigos sobre rendezvous;
% dinâmica relativa;
% aproximação orbital;
% cooperação/não cooperação do alvo;
% evolução temporal da área.

\section{Estimação de pose do alvo e percepção relativa em operações de proximidade}
\label{sec_estimacao_pose_alvo_percepcao_relativa}

Nesta seção é apresentada uma revisão da literatura relacionada à estimação de pose do alvo e à percepção relativa em operações de proximidade orbital. São destacadas as referências que tratam do uso de sensores ópticos e ativos, da reconstrução do estado relativo do alvo e da integração entre navegação relativa, guiamento e controle em missões de \gls{rvdb} com alvos cooperativos e não cooperativos.

Com a evolução das missões de serviço em órbita, remoção ativa de detritos e inspeção de satélites, a percepção relativa passou a ocupar papel central nas operações de proximidade. Nessas missões, tornou-se necessário estimar não apenas a posição relativa entre perseguidor e alvo, mas também sua atitude, velocidade relativa, velocidade angular e, em alguns casos, a própria geometria do alvo. Dessa forma, a literatura passou a explorar diferentes combinações de sensores e algoritmos para a reconstrução do estado relativo durante as fases de aproximação curta.

\cite{volpe2018passive} apresentam um método para determinação da pose e da forma de um satélite não cooperativo e possivelmente desconhecido a partir de imagens obtidas por uma câmera passiva, complementadas por um sensor de distância embarcado na espaçonave perseguidora. Os autores formulam um filtro capaz de estimar posição relativa, atitude, velocidade, velocidade angular e geometria aproximada do alvo com base no rastreamento de características em imagens sucessivas. O trabalho mostra que a navegação relativa pode ser realizada sem depender de marcas cooperativas previamente instaladas no alvo, explorando diretamente informações extraídas da cena observada. Os resultados apresentados indicam desempenho satisfatório na estimação do estado relativo e da forma aproximada do alvo.

Em continuidade a essa linha, \cite{volpe2022gnc} apresentam uma arquitetura integrada de guiamento, navegação e controle para \gls{rvdb} com um alvo não cooperativo em rotação, utilizando câmera monocular passiva. Nesse trabalho, a estimação de pose não é tratada isoladamente, mas integrada a uma arquitetura fechada de navegação e controle. Os autores utilizam extração e rastreamento de características visuais, em conjunto com um filtro \gls{ukf}, para estimar posição relativa, velocidade, atitude, velocidade angular e forma reconstruída do alvo. A partir dessas estimativas, é gerada uma trajetória ótima livre de colisão, validada em co-simulação com imagens renderizadas em malha fechada. Essa referência mostra uma evolução da literatura no sentido da integração entre percepção relativa e decisão de aproximação.

Além das abordagens baseadas em sensores passivos, a literatura recente também apresenta o uso de sensores ativos para contornar limitações associadas à iluminação, à baixa textura superficial do alvo e às distorções provocadas pelo movimento relativo. Nesse contexto, \cite{sun2022tof} apresentam um sistema de estimação de pose de alvos espaciais não cooperativos baseado em câmera \gls{tof}.
Os autores desenvolvem uma arquitetura composta por pré-processamento dos dados, registro quadro a quadro e atualização de modelo global, com o objetivo de reduzir o erro acumulado e o custo computacional. O método é validado por experimentos semi-físicos, apresentando erro translacional mínimo de 8~mm, erro angular de aproximadamente 1$^\circ$, tempo de processamento em torno de 100~ms e frequência de rastreamento de 10~Hz. Esses resultados mostram a viabilidade do uso de sensores ativos compactos em operações de proximidade.

Mais recentemente, \cite{gavilanez2024vision} apresentam uma abordagem para estimação relativa de posição e atitude de espaçonaves não cooperativas baseada em visão estéreo e arquitetura generativa. O método combina extração de pontos característicos, reconstrução de nuvem de pontos por rede generativa adversarial e alinhamento por meio do algoritmo \gls{gicp}.
Os autores buscam melhorar a resolução e a distribuição das nuvens de pontos antes da etapa de registro, com o objetivo de aumentar a precisão final da estimação. Os resultados experimentais em bancada indicam melhora no desempenho da determinação de posição e atitude, sugerindo que técnicas de aprendizado de máquina podem ampliar a robustez de métodos visuais em operações de proximidade.

Paralelamente ao desenvolvimento de métodos específicos, surgiram trabalhos de revisão dedicados à sistematização da navegação relativa baseada em aprendizado profundo. \cite{song2022survey} apresentam uma revisão de métodos de navegação relativa de espaçonaves baseados em \textit{deep learning}, discutindo aplicações em \gls{rvdb}, exploração de pequenos corpos e navegação por terreno, além de bases de dados e desafios de implementação embarcada. Por sua vez, \cite{pauly2023survey} apresentam uma revisão voltada especificamente à estimação monocular de pose de espaçonaves, destacando limitações importantes da literatura recente, como a dificuldade de implantação de modelos computacionalmente intensivos em hardware espacial e a redução de desempenho quando métodos treinados com imagens sintéticas são aplicados a dados mais próximos de condições reais. Essas revisões mostram que, embora o uso de aprendizado profundo tenha ampliado significativamente o repertório técnico da área, ainda permanecem desafios relacionados à confiabilidade, à representatividade dos dados e à transferência entre ambiente simulado e operação real.

De modo geral, a literatura revisada mostra que a estimação de pose do alvo e a percepção relativa passaram a constituir elementos centrais das operações autônomas de proximidade. Enquanto trabalhos anteriores enfatizavam a reconstrução geométrica e a obtenção do estado relativo por sensores ópticos, estudos mais recentes passaram a incorporar integração com arquiteturas de guiamento, navegação e controle, uso de sensores ativos, técnicas baseadas em aprendizado de máquina e preocupações explícitas com robustez computacional e implementação embarcada. Esses aspectos são relevantes para a presente dissertação, pois mostram que a fase de aproximação curta depende não apenas da modelagem da dinâmica relativa, mas também da disponibilidade de informações confiáveis sobre a configuração espacial do alvo ao longo da operação.
% visão computacional;
% sensores;
% estimação de pose;
% navegação relativa;
% percepção do alvo em missões de proximidade.

% \section{Sistemas acoplados espaçonave--manipulador} % modelagem_controle_espaconave_manipulador.tex
% \label{sec_modelagem_controle_espaconave_manipulador}

% Nesta seção é apresentada uma revisão da literatura relacionada à modelagem e ao controle de sistemas acoplados formados por uma espaçonave e um manipulador robótico. São destacadas as referências que tratam da modelagem cinemática e dinâmica desses sistemas, do acoplamento entre a base e o manipulador, das formulações empregadas para a obtenção das equações de movimento e das estratégias de controle aplicadas a manipuladores espaciais.

% %%%%%%%%%%%% modelagem cinematica

% A modelagem cinemática de manipuladores espaciais constitui uma etapa central para descrever o movimento relativo entre os elos, determinar a posição e a orientação do efetuador final e estabelecer as relações entre o espaço das juntas e o espaço cartesiano \cite{fonseca2019kinematics}. No contexto espacial, essa modelagem adquire complexidade adicional, pois o manipulador está acoplado a uma plataforma que, em geral, não pode ser tratada como uma base fixa convencional. Em vez disso, a formulação deve considerar a interação entre o manipulador e a espaçonave, bem como o uso consistente dos referenciais inercial, orbital e fixo ao corpo \cite{fonseca2019kinematics,fonseca2016dynamics}.

% Nesse cenário, abordagens clássicas de cinemática direta e inversa continuam desempenhando papel relevante. A convenção de Denavit--Hartenberg permanece amplamente empregada para a definição sistemática dos referenciais e para a obtenção das transformações homogêneas entre os elos, permitindo calcular a pose do efetuador final a partir das variáveis articulares \cite{fonseca2019kinematics}. Quando a cinemática inversa não admite solução analítica em forma fechada, o problema pode ser formulado como uma otimização, na qual se busca a configuração articular que minimiza o erro entre a pose desejada e a pose obtida pelo manipulador \cite{santos2022uncertainty,santos2023physics}.

% Entretanto, em manipuladores do tipo espaçonave, a formulação cinemática não se limita à descrição geométrica da cadeia serial. Como a movimentação das juntas pode induzir deslocamentos e rotações da plataforma, torna-se necessário compatibilizar a cinemática do manipulador com os movimentos da espaçonave e com os referenciais relevantes ao problema orbital \cite{fonseca2019kinematics,fonseca2016dynamics}. Essa característica é especialmente importante em operações de proximidade, captura e serviço em órbita, nas quais erros de posicionamento e orientação podem comprometer o desempenho da missão \cite{arantes2010far,fonseca2016dynamics}.

% Além das formulações analíticas tradicionais, estudos mais recentes têm explorado estratégias baseadas em aprendizado de máquina para apoiar a resolução da cinemática inversa e o planejamento de movimento. Redes neurais e métodos como máquinas de vetores de suporte vêm sendo empregados para estimar configurações articulares, quantificar incertezas e melhorar a convergência de procedimentos de otimização, inclusive em situações nas quais os efeitos dinâmicos da plataforma e do manipulador são incorporados ao problema \cite{santos2022uncertainty,santos2023physics,santos2022machine}.

% %%%%%%%%%%%%%%%%%%

% Diferentemente dos manipuladores fixos em solo, os manipuladores espaciais operam em sistemas nos quais o movimento das juntas produz reações sobre a base, alterando sua atitude e, em determinadas configurações, também seu movimento translacional.
% Dessa forma, a modelagem e o controle desses sistemas exigem o tratamento explícito do acoplamento dinâmico entre a espaçonave e o manipulador, não sendo suficiente considerar a base como um suporte rigidamente fixo. 
% Conforme estabelecido classicamente por \cite{umetani1989resolved}, \cite{papadopoulos1991nature} e \cite{dubowsky1993kinematics}, a formulação da cinemática e da dinâmica, como o uso da Jacobiana generalizada, deve incorporar a natureza flutuante da base e suas restrições de controle de forma integrada.
% No que se refere aos avanços posteriores na formulação dinâmica desses sistemas, \cite{from2011singularityfree} apresentam equações dinâmicas livres de singularidades para sistemas espaçonave--manipulador, utilizando uma representação mínima baseada em quase-coordenadas em um arcabouço lagrangiano.
% Essa abordagem permite evitar as singularidades inerentes às parametrizações de atitude por ângulos de Euler, resultando em um modelo computacionalmente mais robusto para simulações de malha fechada.

% Os trabalhos de \cite{floresabad2014review} apresentam uma revisão sobre tecnologias de robótica espacial para \textit{on-orbit servicing}, discutindo problemas de cinemática, dinâmica, controle e verificação experimental de sistemas robóticos acoplados. Os autores destacam que a modelagem e o controle de manipuladores espaciais continuam entre os principais desafios das missões autônomas em órbita, especialmente em aplicações de captura, reparo, montagem, reabastecimento e remoção de detritos. Essa revisão é importante por situar os desenvolvimentos teóricos em um contexto operacional mais amplo.

% Os estudos feitos por \cite{fonseca2017attitude} analisam a dinâmica de atitude de uma espaçonave tratada de forma análoga a um manipulador robótico no contexto de missões de \textit{on-orbit servicing}. Os autores formulam o problema com base na dinâmica rotacional em órbita e avaliam o uso de controladores LQR e PID para as formas linearizada e não linear do modelo. As simulações realizadas em MATLAB mostram que ambas as abordagens atendem aos objetivos definidos para operações de \textit{rendezvous} e \textit{berthing}, permitindo ainda comparar esforços de controle e comportamento dinâmico sob diferentes formulações. Essa referência é particularmente relevante para a presente dissertação por enfatizar a necessidade de tratar de forma integrada a atitude da espaçonave e a dinâmica associada ao sistema robótico embarcado.

% Mais recentemente, \cite{rybus2017nmpc} apresentam um sistema de controle para manipuladores espaciais \textit{free-floating} baseado em controle preditivo não linear. A arquitetura proposta combina planejamento de trajetória com um controlador NMPC, ambos formulados de modo a considerar explicitamente a natureza flutuante do sistema espaçonave--manipulador. O método é avaliado por simulações numéricas em um caso planar simplificado, incluindo perturbações e incertezas paramétricas. Os resultados apresentados indicam desempenho superior ao de estratégias baseadas na inversa da Jacobiana dinâmica e no \textit{Modified Simple Adaptive Control}. Essa referência mostra uma tendência mais recente da literatura em direção a estratégias de controle mais integradas, capazes de lidar com rastreamento, restrições dinâmicas e robustez a incertezas.

% De modo geral, a literatura revisada mostra que a modelagem e o controle de sistemas acoplados espaçonave--manipulador evoluíram de formulações cinemáticas fundamentais para representações dinâmicas mais robustas e, posteriormente, para arquiteturas de controle mais sofisticadas. Os trabalhos revisados mostram que o tratamento do acoplamento entre a espaçonave e o manipulador constitui aspecto central do problema, sendo indispensável para a descrição adequada do comportamento do sistema em operações orbitais. Esses aspectos são relevantes para a presente dissertação, pois reforçam a necessidade de formular de modo coerente a modelagem e o controle do sistema acoplado nas fases de aproximação e interação com o alvo.

\section{Sistemas acoplados espaçonave--manipulador} % modelagem_controle_espaconave_manipulador.tex
\label{sec_modelagem_controle_espaconave_manipulador}

Nesta seção é apresentada uma revisão da literatura relacionada à modelagem e ao controle de sistemas acoplados formados por uma espaçonave e um manipulador robótico. São destacadas as referências que tratam da modelagem cinemática e dinâmica desses sistemas, do acoplamento entre a base e o manipulador, das formulações empregadas para a obtenção das equações de movimento e das estratégias de controle aplicadas a manipuladores espaciais.

A modelagem cinemática de manipuladores espaciais constitui uma etapa central para descrever o movimento relativo entre os elos, determinar a posição e a orientação do efetuador final e estabelecer as relações entre o espaço das juntas e o espaço cartesiano \cite{fonseca2019kinematics}. No contexto espacial, essa modelagem adquire complexidade adicional, pois o manipulador está acoplado a uma plataforma que, em geral, não pode ser tratada como uma base fixa convencional. Em vez disso, a formulação deve considerar a interação entre o manipulador e a espaçonave, bem como o uso consistente dos referenciais inercial, orbital e fixo ao corpo \cite{fonseca2019kinematics,fonseca2016dynamics}.

Nesse cenário, abordagens clássicas de cinemática direta e inversa continuam desempenhando papel relevante. A convenção de Denavit--Hartenberg permanece amplamente empregada para a definição sistemática dos referenciais e para a obtenção das transformações homogêneas entre os elos, permitindo calcular a pose do efetuador final a partir das variáveis articulares \cite{fonseca2019kinematics}. Quando a cinemática inversa não admite solução analítica em forma fechada, o problema pode ser formulado como uma otimização, na qual se busca a configuração articular que minimiza o erro entre a pose desejada e a pose obtida pelo manipulador \cite{santos2022uncertainty,santos2023physics}.

Entretanto, em manipuladores do tipo espaçonave, a formulação cinemática não se limita à descrição geométrica da cadeia serial. Como a movimentação das juntas pode induzir deslocamentos e rotações da plataforma, torna-se necessário compatibilizar a cinemática do manipulador com os movimentos da espaçonave e com os referenciais relevantes ao problema orbital \cite{fonseca2019kinematics,fonseca2016dynamics}. Essa característica é especialmente importante em operações de proximidade, captura e serviço em órbita, nas quais erros de posicionamento e orientação podem comprometer o desempenho da missão \cite{arantes2010far,fonseca2016dynamics}.

Além das formulações analíticas tradicionais, estudos mais recentes têm explorado estratégias baseadas em aprendizado de máquina para apoiar a resolução da cinemática inversa e o planejamento de movimento. Redes neurais e métodos como máquinas de vetores de suporte vêm sendo empregados para estimar configurações articulares, quantificar incertezas e melhorar a convergência de procedimentos de otimização, inclusive em situações nas quais os efeitos dinâmicos da plataforma e do manipulador são incorporados ao problema \cite{santos2022uncertainty,santos2023physics,santos2022machine}.

Do ponto de vista dinâmico, os manipuladores espaciais diferem dos manipuladores fixos em solo, pois operam em sistemas nos quais o movimento das juntas produz reações sobre a base, alterando sua atitude e, em determinadas configurações, também seu movimento translacional. Dessa forma, a modelagem e o controle desses sistemas exigem o tratamento explícito do acoplamento entre a espaçonave e o manipulador, não sendo suficiente considerar a base como um suporte rigidamente fixo. Conforme estabelecido classicamente por \cite{umetani1989resolved}, \cite{papadopoulos1991nature} e \cite{dubowsky1993kinematics}, a formulação da cinemática e da dinâmica, incluindo o uso da Jacobiana generalizada, deve incorporar, de forma integrada, a natureza flutuante da base e as restrições de controle associadas a esse tipo de sistema.

No que se refere aos avanços posteriores na formulação dinâmica desses sistemas, \cite{from2011singularityfree} apresentam equações dinâmicas livres de singularidades para sistemas espaçonave--manipulador, utilizando uma representação mínima baseada em quase-coordenadas em um arcabouço lagrangiano. Essa abordagem permite evitar as singularidades inerentes às parametrizações de atitude por ângulos de Euler, resultando em um modelo computacionalmente mais robusto para simulações em malha fechada.

Os trabalhos de \cite{floresabad2014review} apresentam uma revisão sobre tecnologias de robótica espacial para \textit{on-orbit servicing}, discutindo problemas de cinemática, dinâmica, controle e verificação experimental de sistemas robóticos acoplados. Os autores destacam que a modelagem e o controle de manipuladores espaciais continuam entre os principais desafios das missões autônomas em órbita, especialmente em aplicações de captura, reparo, montagem, reabastecimento e remoção de detritos. Essa revisão é importante por situar os desenvolvimentos teóricos em um contexto operacional mais amplo.

Os estudos de \cite{fonseca2017attitude} analisam a dinâmica de atitude de uma espaçonave tratada de forma análoga a um manipulador robótico no contexto de missões de \textit{on-orbit servicing}. Os autores formulam o problema com base na dinâmica rotacional em órbita e avaliam o uso de controladores \gls{lqr} e \gls{pid} para as formas linearizada e não linear do modelo. As simulações realizadas em MATLAB mostram que ambas as abordagens atendem aos objetivos definidos para operações de \textit{rendezvous} e \textit{berthing}, permitindo ainda comparar esforços de controle e comportamento dinâmico sob diferentes formulações. Essa referência é particularmente relevante para a presente dissertação, pois enfatiza a necessidade de tratar de forma integrada a atitude da espaçonave e a dinâmica associada ao sistema robótico embarcado.

Mais recentemente, \cite{rybus2017nmpc} apresentam um sistema de controle para manipuladores espaciais \textit{free-floating} baseado em controle preditivo não linear. A arquitetura proposta combina planejamento de trajetória com um controlador \gls{nmpc}, ambos formulados de modo a considerar explicitamente a natureza flutuante do sistema espaçonave--manipulador. O método é avaliado por simulações numéricas em um caso planar simplificado, incluindo perturbações e incertezas paramétricas. Os resultados apresentados indicam desempenho superior ao de estratégias baseadas na inversa da Jacobiana dinâmica e no \textit{Modified Simple Adaptive Control}. Essa referência mostra uma tendência mais recente da literatura em direção a estratégias de controle mais integradas, capazes de lidar com rastreamento, restrições dinâmicas e robustez a incertezas.

De modo geral, a literatura revisada mostra que a modelagem e o controle de sistemas acoplados espaçonave--manipulador evoluíram de formulações cinemáticas fundamentais para representações dinâmicas mais robustas e, posteriormente, para arquiteturas de controle mais sofisticadas. Os trabalhos revisados mostram que o tratamento do acoplamento entre a espaçonave e o manipulador constitui um aspecto central do problema, sendo indispensável para a descrição adequada do comportamento do sistema em operações orbitais. Esses aspectos são relevantes para a presente dissertação, pois reforçam a necessidade de formular, de modo coerente, a modelagem e o controle do sistema acoplado nas fases de aproximação e interação com o alvo.
% modelagem cinemática e dinâmica;
% acoplamento entre base e manipulador;
% formulações lagrangianas/newton-euler;
% controle aplicado ao sistema acoplado.

\section{Controle com reação nula e manipuladores \textit{free-floating}}
\label{sec_reaction_null_free_floating}

Nesta seção é apresentada uma revisão da literatura relacionada ao controle de manipuladores espaciais \textit{free-floating} com redução ou anulação das reações transmitidas à base. São destacadas as referências que tratam do conceito de \textit{reaction null space}, do controle \textit{reactionless}, da captura de alvos com mínima perturbação da base e de estratégias adaptativas aplicadas a sistemas espaciais acoplados.

Entre as contribuições específicas da literatura de robótica espacial, destaca-se a linha de pesquisa voltada ao controle de manipuladores \textit{free-floating} com redução ou anulação das reações dinâmicas impostas à base. Nessa classe de problemas, o objetivo não é apenas garantir o rastreamento da trajetória do efetuador terminal, mas fazê-lo de modo a preservar a atitude da espaçonave ou a integridade dinâmica da estrutura que suporta o manipulador. Essa questão é particularmente importante em missões de proximidade, captura e serviço em órbita, nas quais perturbações excessivas da base podem comprometer a segurança da operação e aumentar o consumo de atuadores auxiliares.

Um dos trabalhos mais importantes nessa linha é apresentado em \cite{nenchev1999reaction}, onde é proposta uma lei de controle composta para o rastreamento da trajetória do efetuador terminal em um sistema manipulador montado sobre estrutura flexível, de forma a não induzir perturbações na base. Os autores formulam o problema com base no conceito de \textit{reaction null space}, introduzido para tratar a interação dinâmica em robôs \textit{free-floating} e sistemas de base móvel em geral.
Além da reação nula, o trabalho inclui a supressão de vibração da base e apresenta validação por simulação e por meio do banco experimental, desenvolvido na \gls{trep}.
Essa referência tornou-se importante por consolidar o conceito de \textit{reaction null space} como ferramenta de projeto para desacoplar, ao menos parcialmente, o movimento útil do manipulador das perturbações transmitidas à estrutura de suporte.

\cite{piersigilli2010reactionless} investigam o problema da captura \textit{reactionless} de um satélite por um manipulador de dois graus de liberdade. Nesse trabalho, os autores tratam explicitamente a captura sem reação de um alvo por um manipulador espacial, reforçando a aplicação prática desse tipo de estratégia em operações de contato e apreensão. Essa referência é relevante porque desloca a discussão de reação nula do campo conceitual para uma tarefa operacional crítica, na qual impulsos e perturbações na base tornam-se ainda mais sensíveis.

Para ampliar essa discussão, \cite{nguyenhuynh2013adaptive} consideram o problema de movimento \textit{reactionless} adaptativo e identificação de parâmetros após a captura de detritos espaciais. Os autores mostram que a captura de um alvo altera a dinâmica do sistema e introduz incertezas paramétricas adicionais, tornando insuficientes estratégias baseadas em modelos exatamente conhecidos. Dessa forma, a combinação entre redução de reação e identificação paramétrica passa a ocupar papel importante na literatura, aproximando o problema de cenários mais realistas de serviço em órbita e pós-captura.

\cite{xu2016adaptive} apresentam um controle adaptativo de movimento com reação nula para manipuladores espaciais \textit{free-floating} com incertezas cinemáticas e dinâmicas. Os autores destacam que uma das principais dificuldades para a construção de uma lei adaptativa baseada em \textit{reaction null space} está na obtenção de uma expressão linear em relação aos parâmetros incertos. A partir dessa reformulação, é proposto um controlador em nível de velocidade capaz de regular a atitude da espaçonave e realizar o rastreamento da trajetória do efetuador terminal. O método é demonstrado por simulações numéricas em um manipulador planar de três graus de liberdade. Essa referência mostra um avanço da literatura ao levar a ideia de reação nula para um cenário com incertezas mais próximas das encontradas em missões reais.

Mais recentemente, \cite{ye2019adaptive} apresentam uma estratégia adaptativa de planejamento e controle baseada em \textit{Adaptive Reaction Null Space} para robôs espaciais \textit{free-floating} em contato com alvos desconhecidos. O método combina um procedimento de planejamento adaptativo via \gls{vffrls} com uma estratégia robusta de controle baseada em \gls{ssadeelm}, buscando compensar distúrbios desconhecidos e incertezas do modelo sem exigir conhecimento preciso dos atributos do alvo.
Os autores apresentam comparações por simulação e experimento com métodos existentes, indicando efetividade e superioridade da estratégia proposta no cenário estudado. Essa referência mostra a evolução da literatura em direção a abordagens adaptativas voltadas a incertezas operacionais e interação com alvos reais.

De modo geral, a literatura revisada mostra que o problema do controle de manipuladores espaciais não se resume ao rastreamento geométrico de trajetórias, envolvendo também a gestão das reações dinâmicas impostas à base. Os trabalhos analisados mostram uma evolução que vai das formulações conceituais de \textit{reaction null space} e experimentos em estruturas flexíveis até estratégias adaptativas voltadas à captura, ao pós-captura e ao contato com alvos desconhecidos.
Esses aspectos são relevantes para a presente dissertação, pois reforçam que, em sistemas acoplados espaçonave--manipulador, o desempenho do braço robótico e o comportamento dinâmico da espaçonave não podem ser tratados de forma dissociada, especialmente nas fases de aproximação final e interação física com o alvo.
% reaction null space;
% zero reaction;
% free-floating space manipulators;
% controle sem perturbar a base;
% artigos que seu orientador explicitamente pediu.

\section{Planejamento, autonomia e técnicas recentes aplicadas a operações de proximidade}
\label{sec_planejamento_autonomia_tecnicas_recentes}

Nesta seção é apresentada uma revisão da literatura relacionada ao planejamento, à autonomia e às técnicas recentes aplicadas a operações de proximidade orbital. São destacadas as referências que tratam do planejamento de trajetória, do guiamento sob restrições, da prevenção de colisão, da tomada de decisão autônoma e da incorporação de técnicas recentes, incluindo métodos baseados em aprendizado de máquina, em missões de \textit{rendezvous}, \textit{docking} e \textit{on-orbit servicing}.

Nos desenvolvimentos mais recentes da literatura, as operações de proximidade passaram a ser tratadas não apenas como problemas de rastreamento de trajetória ou de estabilização local, mas também como tarefas integradas de percepção, planejamento, tomada de decisão e controle sob restrições. Essa evolução está associada à crescente complexidade das missões de \textit{on-orbit servicing}, remoção ativa de detritos e \textit{rendezvous} com alvos não cooperativos, nas quais a espaçonave perseguidora deve operar em ambientes geometricamente restritos, com incertezas de estado, risco de colisão e necessidade de maior autonomia.

\cite{alizadeh2024survey} apresentam uma revisão sobre robótica espacial para \textit{on-orbit servicing}, destacando a interação entre estimação do estado do alvo, planejamento de movimento, controle em malha fechada e técnicas de aprendizado de máquina. Os autores mostram que o avanço recente da área depende cada vez mais da integração entre essas camadas funcionais, e não apenas do desenvolvimento isolado de subsistemas. Essa referência é importante por situar o problema do planejamento e da autonomia em um contexto mais amplo de operação robótica espacial.

\cite{dong2022safe} apresentam uma estrutura de controle preditivo voltada ao \textit{rendezvous} e \textit{docking} autônomos com alvo cambaleante em múltiplos estágios. Os autores consideram restrições práticas como saturação de atuadores, limites de velocidade e prevenção de colisão em condições adequadas ao acoplamento final. A formulação é construída de modo a garantir segurança e estabilidade, combinando etapas de aproximação global e refinamento terminal. Os resultados apresentados mostram que a abordagem conduz a espaçonave perseguidora a regiões compatíveis com o acoplamento, preservando requisitos de segurança operacional. Essa referência mostra uma tendência importante da literatura recente em direção a arquiteturas de planejamento e controle formuladas explicitamente sob restrições.

O problema de manobras de proximidade com controle preditivo é estudado por \cite{wang2023close} para \textit{rendezvous} e \textit{docking} com cliente rotativo ou não rotativo, modelando explicitamente a \textit{keep-out zone} do alvo por meio de uma elipsoide expandida. Os autores buscam melhorar simultaneamente o desempenho e a eficiência computacional do controlador, aspecto importante para aplicações em tempo real. Essa referência é relevante porque aproxima a formulação matemática das condições operacionais encontradas em missões reais, nas quais não basta convergir à vizinhança do alvo, sendo necessário respeitar zonas seguras, evitar colisões e manter viabilidade computacional ao longo de toda a aproximação.

Mais recentemente, \cite{madonna2025guidance} apresentam uma arquitetura de guiamento e controle para \textit{rendezvous} a curta distância e aproximação final de uma espaçonave perseguidora em direção a um alvo não cooperativo. O estudo combina planejamento de trajetória, lógica de guiamento e tratamento explícito da geometria da aproximação final, com foco em operações de curta distância. Essa referência mostra uma tendência mais recente da literatura, na qual a fase terminal da operação passa a ser tratada como uma etapa com requisitos geométricos, cinemáticos e operacionais próprios, e não apenas como um prolongamento da dinâmica relativa linearizada.

Paralelamente às abordagens baseadas em otimização e controle preditivo, a literatura recente também passou a explorar técnicas de aprendizado por reforço para ampliar a autonomia decisória em cenários com obstáculos, incertezas e necessidade de resposta rápida. \cite{sharma2024safety} apresentam uma metodologia de guiamento para \textit{rendezvous} baseada em \textit{safety reinforcement learning}, formulando o problema como um processo de decisão de Markov em situações que envolvem prevenção de colisão. Os autores mostram que a estratégia proposta é capaz de realizar desvio de obstáculos e atingir \textit{rendezvous} de alta precisão, ao mesmo tempo em que busca lidar com uma limitação recorrente dos métodos tradicionais de aprendizado por reforço, associada à dificuldade de impor restrições de segurança durante a exploração. Essa referência mostra que a autonomia em missões de proximidade passou a envolver não apenas automação do controle, mas também maior delegação do processo decisório.

No contexto da robótica espacial acoplada, \cite{alali2024reinforcement} investigam o planejamento de trajetória para um manipulador espacial \textit{free-floating} de seis graus de liberdade, com foco em prevenção de colisão e desvio de obstáculos por meio de aprendizado por reforço. Os autores consideram explicitamente o efeito do acoplamento dinâmico entre a espaçonave e o manipulador, aspecto particularmente importante nesse tipo de sistema. Essa referência amplia a aplicação do aprendizado por reforço para além da aproximação translacional entre espaçonaves, alcançando também o problema de planejamento de movimento do manipulador em ambiente orbital.

De modo geral, a literatura revisada mostra que as pesquisas mais recentes caminham em direção a sistemas mais autônomos, integrados e orientados por restrições reais de missão. Enquanto trabalhos anteriores concentravam-se principalmente na formulação da dinâmica relativa e no controle de seguimento, as contribuições recentes passaram a enfatizar arquiteturas mais completas, que combinam percepção, planejamento, controle, prevenção de colisão e tomada de decisão assistida por técnicas de otimização e aprendizado. Esses aspectos são relevantes para a presente dissertação, pois mostram que a fase de aproximação e interação com o alvo deve ser compreendida não apenas como um problema dinâmico, mas também como parte de um processo operacional mais amplo, no qual segurança, autonomia e integração entre subsistemas assumem papel central.
% planejamento de trajetória;
% guiamento;
% autonomia;
% IA integrada;
% prevenção de colisão;
% tomada de decisão;
% trabalhos recentes e, quando houver, validação experimental.

\section{Síntese crítica da literatura e posicionamento da dissertação}
\label{sec_sintese_critica_posicionamento_dissertacao}

A revisão da literatura apresentada ao longo deste capítulo evidencia que os estudos sobre operações de \textit{rendezvous} e \textit{berthing} evoluíram de formulações mais clássicas de dinâmica relativa para abordagens mais amplas, envolvendo percepção relativa, controle de atitude, planejamento de trajetória, autonomia e interação robótica. As referências revisadas indicam que a literatura avançou tanto na formulação matemática do movimento relativo entre espaçonaves quanto na modelagem e no controle de sistemas robóticos espaciais acoplados.

No campo da dinâmica relativa, os trabalhos clássicos estabeleceram as bases matemáticas para a descrição do movimento entre perseguidor e alvo. Em estudos posteriores, o problema passou a incorporar restrições operacionais, requisitos de segurança e condições associadas à aproximação de alvos cooperativos e não cooperativos. Na área de percepção relativa, a literatura passou a apresentar métodos voltados à estimação de posição, atitude, velocidade relativa e, em alguns casos, da própria geometria do alvo, utilizando sensores passivos e ativos, além de técnicas mais recentes baseadas em aprendizado de máquina.

Além disso, a literatura revisada mostra que as operações de \textit{rendezvous} e \textit{berthing} vêm sendo tratadas em contextos operacionais mais amplos, envolvendo diferentes níveis de cooperação do alvo, troca de informações entre sistemas embarcados, requisitos de segurança e maior autonomia decisória ao longo da missão.
Embora nem todas essas abordagens constituam o núcleo da metodologia adotada nesta dissertação, sua inclusão na revisão permite situar o estado da arte da área e delimitar com maior clareza o recorte adotado no presente trabalho.

No que se refere à robótica espacial, a literatura revisada evidencia que os sistemas espaçonave--manipulador devem ser tratados como sistemas dinamicamente acoplados, nos quais o movimento das juntas produz reações sobre a base e altera a dinâmica global do sistema. As referências analisadas indicam ainda que estratégias como \textit{reaction null space}, controle \textit{reactionless} e abordagens adaptativas ampliaram a capacidade de tratar problemas de captura, pós-captura e interação com alvos desconhecidos, preservando, tanto quanto possível, o comportamento dinâmico da base.

Entretanto, a revisão também mostra que parte significativa da literatura trata esses problemas de forma parcial. Em muitos trabalhos, a dinâmica relativa orbital é formulada separadamente da dinâmica de atitude da espaçonave perseguidora. Em outros, a modelagem do manipulador robótico é desenvolvida sem integração explícita com o contexto de aproximação orbital de curta distância. Há ainda referências voltadas à percepção relativa, à estimação de pose ou ao planejamento autônomo que, embora importantes, não tratam de forma integrada os efeitos dinâmicos do sistema completo quando a espaçonave perseguidora transporta ou opera um manipulador robótico em ambiente orbital.

Dessa forma, a principal lacuna observada na literatura revisada não está na ausência de estudos sobre os diferentes elementos do problema, mas na limitada integração entre eles dentro de uma mesma formulação de modelagem e análise. Em particular, permanece relevante o estudo de formulações que articulem a dinâmica relativa entre perseguidor e alvo, a dinâmica de atitude da espaçonave perseguidora no referencial orbital e a dinâmica acoplada de um manipulador robótico embarcado, com vistas à compreensão do comportamento global do sistema durante operações de aproximação e interação com o alvo.

É nesse contexto que se insere a presente dissertação. O trabalho concentra-se na modelagem integrada do sistema formado por uma espaçonave perseguidora equipada com manipulador robótico, considerando o contexto de operações de \textit{rendezvous} e \textit{berthing} em ambiente orbital. O foco recai sobre a articulação entre dinâmica relativa, dinâmica de atitude e modelagem do manipulador, de modo a representar de forma consistente os acoplamentos envolvidos nas fases de aproximação e nas condições que antecedem a interação com o alvo.

Assim, a contribuição principal desta dissertação não está na proposição de uma arquitetura completa de autonomia, nem no desenvolvimento de um sistema embarcado de percepção relativa ou de tomada de decisão em tempo real. O objetivo principal está na construção e na análise de um modelo dinâmico coerente com o problema físico estudado, capaz de servir de base para investigações posteriores em controle, planejamento e operação robótica espacial. Dessa forma, o presente trabalho adota um escopo mais delimitado, concentrando-se na formulação matemática e na análise do sistema acoplado em ambiente orbital.

De modo geral, a revisão da literatura evidencia que, embora a área tenha avançado de forma significativa em dinâmica relativa, percepção, controle e robótica espacial, ainda permanece relevante o desenvolvimento de abordagens que privilegiem a integração entre esses domínios. Nesse sentido, a presente dissertação busca contribuir ao tratar de forma articulada a dinâmica orbital relativa, a atitude da espaçonave perseguidora e a dinâmica do manipulador robótico, definindo, assim, seu posicionamento no conjunto das pesquisas dedicadas às operações de \textit{rendezvous} e \textit{berthing} com apoio de sistemas robóticos espaciais.
% síntese final;
% lacunas reais demonstradas pela revisão;
% diferenciais do seu trabalho;
% delimitação do que sua dissertação faz e do que não faz.
